\ifdefined\gkver\else\def\gkver{student}\fi
\documentclass[\gkver]{gaokaozhenti}
\begin{document}

\chapter{2017年海南理综}
%% paperType: 理综物理部分

\begin{enumerate}
\item
%% number: 1
%% paperName: 2017年海南理综第 1 题 4 分
%% typeId: 1
%% type: 单选题
%% chapter: 运动和力
%% point: 动量守恒定律
%% method: 无
%% score: 4
%% degree: 850
%% duplicateId: 0
%% body:
光滑水平桌面上有 P、Q 两个物块，Q 的质量是 P 的 $n$ 倍。将一轻弹簧置于 P、Q 之间，用外力缓慢压 P、Q。撤去外力后，P、Q 开始运动，P 和 Q 的动量大小的比值为\xzanswer{D}
\fourchoices[answer=D]
{$n^{2}$}
{$n$}
{$\frac{1}{n}$}
{$1$}
%% answer: D
%% memo:
\memoanswer{
撤去外力后，系统不受外力，所以总动量守恒。设 P 的动量方向为正方向，则有 $p_{P}-p_{Q}=0$，故 $p_{P}=p_{Q}$，P 和 Q 的动量大小之比为 $1$；故 D 正确，ABC 错误。
}

\item
%% number: 2
%% paperName: 2017年海南理综第 2 题 4 分
%% typeId: 1
%% type: 单选题
%% chapter: 电场
%% point: 电场线
%% method: 无
%% score: 4
%% degree: 750
%% duplicateId: 0
%% body:
关于静电场的电场线，下列说法正确的是\xzanswer{C}
\fourchoices[answer=C]
{电场强度较大的地方电场线一定较疏}
{沿电场线方向，电场强度一定越来越小}
{沿电场线方向，电势一定越来越低}
{电场线一定是带电粒子在电场中运动的轨迹}
%% answer: C
%% memo:
\memoanswer{
A、电场线的疏密表示场强的强弱，那么电场强度较大的地方电场线一定较密，故 A 错误；BC、沿着电场线的方向，电势会降低，因此沿电场线方向电势越来越低，但电场线不一定越来越疏，则场强不一定越来越小，故 B 错误，C 正确；D、电场线不一定与带电粒子的轨迹重合，只有电场线是直线，带电粒子的初速度为零或初速度方向与电场线方向在同一条直线上时电场线才与带电粒子的轨迹重合，故 D 错误。故选：C。
}

\item
%% number: 3
%% paperName: 2017年海南理综第 3 题 4 分
%% typeId: 1
%% type: 单选题
%% chapter: 运动和力
%% point: 牛顿定律的应用
%% method: 无
%% score: 4
%% degree: 750
%% duplicateId: 0
%% body:
汽车紧急刹车后，停止运动的车轮在水平地面上滑动直至停止，在地面上留下的痕迹称为刹车线。由刹车线的长短可知汽车刹车前的速度。已知汽车轮胎与地面之间的动摩擦因数为 0.80，测得刹车线长 $25\Um$。汽车在刹车前的瞬间的速度大小为（重力加速度 $g$ 取 $10\Umsq$）\xzanswer{B}
\fourchoices[answer=B]
{$10\Ums$}
{$20\Ums$}
{$30\Ums$}
{$40\Ums$}
%% answer: B
%% memo:
\memoanswer{
刹车后汽车的合外力为摩擦力 $f=\mu mg$，加速度 $a=\frac{f}{m}=\mu g=8\Umsq$；又有刹车线长 $25\Um$，故可由匀变速直线运动规律得到汽车在刹车前的瞬间的速度大小 $v=\sqrt{2as}=\sqrt{2\times8\times25}\Ums=20\Ums$；故 ACD 错误，B 正确；故选：B。
}

\item
%% number: 4
%% paperName: 2017年海南理综第 4 题 4 分
%% typeId: 1
%% type: 单选题
%% chapter: 电场
%% point: 带电粒子的平衡
%% method: 无
%% score: 4
%% degree: 550
%% duplicateId: 0
%% body:
如图，平行板电容器的两极板竖直放置并分别与电源的正负极相连，一带电小球经绝缘轻绳悬挂于两极板之间，处于静止状态。现保持右极板不动，将左极板向左缓慢移动。关于小球所受的电场力大小 $F$ 和绳子的拉力大小 $T$，下列判断正确的是\xzanswer{D}
\onepicture[w=4.2cm]{figs/04.png}
\fourchoices[answer=D]
{$F$ 逐渐减小，$T$ 逐渐减小}
{$F$ 逐渐增大，$T$ 逐渐减小}
{$F$ 逐渐减小，$T$ 逐渐增大}
{$F$ 逐渐增大，$T$ 逐渐增大}
%% answer: D
%% memo:
\memoanswer{
电容器与电源相连，所以两端间电势差不变，将左极板向左缓慢移动过程中，两板间距离减小，则由 $U=Ed$ 可知，电场强度 $E$ 增大；电场力 $F=Eq$ 增大；小球处于平衡状态，受重力、拉力与电场力的作用而处于平衡，故拉力与电场力和重力的合力大小相等，方向相反；根据平行四边形定则可知，$T=\sqrt{F^{2}+(mg)^{2}}$；由于重力不变，电场力变大，故拉力增大。故 D 正确，ABC 错误。故选：D。

\onepicture[w=2.8cm, align=right]{figs/04_an.png}
}

\item
%% number: 5
%% paperName: 2017年海南理综第 5 题 4 分
%% typeId: 1
%% type: 单选题
%% chapter: 曲线运动
%% point: 平抛运动
%% method: 无
%% score: 4
%% degree: 550
%% duplicateId: 0
%% body:
已知地球质量为月球质量的 81 倍，地球半径约为月球半径的 4 倍。若在月球和地球表面同样高度处，以相同的初速度水平抛出物体，抛出点与落地点间的水平距离分别为 $s_{\text{月}}$ 和 $s_{\text{地}}$，则 $s_{\text{月}}:s_{\text{地}}$ 约为\xzanswer{A}
\fourchoices[answer=A]
{$9:4$}
{$6:1$}
{$3:2$}
{$1:1$}
%% answer: A
%% memo:
\memoanswer{
由 $G\frac{mM}{R^{2}}=mg$ 可得在月球和地球表面的重力加速度之比为 $\frac{g_{\text{月}}}{g_{\text{地}}}=\frac{M_{\text{月}}R_{\text{地}}^{2}}{M_{\text{地}}R_{\text{月}}^{2}}=\frac{4^{2}}{81}=\frac{16}{81}$①，由平抛运动规律有 $s=vt$，$h=\frac{1}{2}gt^{2}$②，联立①②解得 $\frac{s_{\text{月}}}{s_{\text{地}}}=\sqrt{\frac{g_{\text{地}}}{g_{\text{月}}}}=\frac{9}{4}$，故 A 正确。
}

\item
%% number: 6
%% paperName: 2017年海南理综第 6 题 4 分
%% typeId: 1
%% type: 单选题
%% chapter: 功和能
%% point: 动能定理
%% method: 无
%% score: 4
%% degree: 550
%% duplicateId: 0
%% body:
将一小球竖直向上抛出，小球在运动过程中所受到的空气阻力不可忽略。a 为小球运动轨迹上的一点，小球上升和下降经过 a 点时的动能分别为 $E_{k1}$ 和 $E_{k2}$。从抛出开始到小球第一次经过 a 点时重力所做的功为 $W_{1}$，从抛出开始到小球第二次经过 a 点时重力所做的功为 $W_{2}$。下列选项正确的是\xzanswer{B}
\fourchoices[answer=B]
{$E_{k1}=E_{k2}$，$W_{1}=W_{2}$}
{$E_{k1}>E_{k2}$，$W_{1}=W_{2}$}
{$E_{k1}<E_{k2}$，$W_{1}<W_{2}$}
{$E_{k1}>E_{k2}$，$W_{1}<W_{2}$}
%% answer: B
%% memo:
\memoanswer{
由重力做功与路径无关可知 $W_{1}=W_{2}$，在小球运动过程中有空气阻力做负功，由动能定理可知，$E_{k1}>E_{k2}$，故 B 正确。
}

\item
%% number: 7
%% paperName: 2017年海南理综第 7 题 5 分
%% typeId: 2
%% type: 多选题
%% chapter: 光学
%% point: 光电效应
%% method: 无
%% score: 5
%% degree: 650
%% duplicateId: 0
%% body:
三束单色光 1、2 和 3 的波长分别为 $\lambda_{1}$、$\lambda_{2}$ 和 $\lambda_{3}$（$\lambda_{1}>\lambda_{2}>\lambda_{3}$）。分别用着三束光照射同一种金属。已知用光束 2 照射时，恰能产生光电子。下列说法正确的是\xzanswer{AC}
\fourchoices[answer=AC]
{用光束 1 照射时，不能产生光电子}
{用光束 3 照射时，不能产生光电子}
{用光束 2 照射时，光越强，单位时间内产生的光电子数目越多}
{用光束 2 照射时，光越强，产生的光电子的最大初动能越大}
%% answer: AC
%% memo:
\memoanswer{
由题述，用光束 2 照射时，恰能产生光电子，可知 $\lambda_{2}$ 为金属的极限波长，又 $\lambda_{1}>\lambda_{2}>\lambda_{3}$，用光束 1 照射时，不能产生光电子，用光束 3 照射时，一定能产生光电子，故 A 正确 B 错误；用光束 2 照射时，光越强，单位时间内照射时产生的光子越多，即产生的光电子数目越多，故 C 正确；由光电效应方程 $E_{k}=h\nu-W_{0}$ 可知，对同一金属，光电子的最大初动能只与入射光的频率有关，故用光束 2 照射时，光越强，产生的光电子数目越多，但其最大初动能不变，故 D 错误。
}

\item
%% number: 8
%% paperName: 2017年海南理综第 8 题 5 分
%% typeId: 2
%% type: 多选题
%% chapter: 交流电
%% point: LC振荡回路
%% method: 无
%% score: 5
%% degree: 650
%% duplicateId: 0
%% body:
如图，电阻 $R$、电容 $C$ 和电感 $L$ 并联后，接入输出电压有效值、频率可调的交流电源。当电路中交流电的频率为 $f$ 时，通过 $R$、$C$ 和 $L$ 的电流有效值恰好相等。若将频率降低为 $\frac{1}{2}f$，分别用 $I_{1}$、$I_{2}$ 和 $I_{3}$ 表示此时通过 $R$、$C$ 和 $L$ 的电流有效值，则\xzanswer{BC}
\onepicture[w=6cm]{figs/08.png}
\fourchoices[answer=BC]
{$I_{1}>I_{3}$}
{$I_{1}>I_{2}$}
{$I_{3}>I_{2}$}
{$I_{2}=I_{3}$}
%% answer: BC
%% memo:
\memoanswer{
由电容器“通高频、阻低频”，电感线圈“通低频、阻高频”可知，若将频率降低为 $\frac{1}{2}f$，则感抗减小，容抗增大，定值电阻阻值不变，通过 $R$ 的电流有效值不变，通过 $C$ 的电流有效值减小，通过 $L$ 的电流有效值增大，即 $I_{3}>I_{1}>I_{2}$，故 BC 正确。
}

\item
%% number: 9
%% paperName: 2017年海南理综第 9 题 5 分
%% typeId: 2
%% type: 多选题
%% chapter: 运动和力
%% point: 连接体
%% method: 无
%% score: 5
%% degree: 450
%% duplicateId: 0
%% body:
如图，水平地面上有三个靠在一起的物块 P、Q 和 R，质量分别为 $m$、$2m$ 和 $3m$，物块与地面间的动摩擦因数都为 $\mu$。用大小为 $F$ 的水平外力推动物块 P，记 R 和 Q 之间相互作用力与 Q 与 P 之间相互作用力大小之比为 $k$。下列判断正确的是\xzanswer{BD}
\onepicture[w=6.5cm]{\includesvg{figs/09}}
\fourchoices[answer=BD]
{若 $\mu\neq0$，则 $k=\frac{5}{6}$}
{若 $\mu\neq0$，则 $k=\frac{3}{5}$}
{若 $\mu=0$，则 $k=\frac{1}{2}$}
{若 $\mu=0$，则 $k=\frac{3}{5}$}
%% answer: BD
%% memo:
\memoanswer{
若 $\mu\neq0$，对系统整体，由牛顿第二定律得 $F-\mu\cdot6mg=6ma$；隔离 R，受力分析，由牛顿第二定律得 $F_{QR}-\mu\cdot3mg=3ma$；隔离 QR，受力分析，由牛顿第二定律得 $F_{PQ}-\mu\cdot5mg=5ma$，联立解得 $k=\frac{3}{5}$，故 B 正确 A 错误；若 $\mu=0$，对整体，由牛顿第二定律得 $F=6ma$；隔离 R，受力分析，由牛顿第二定律得 $F_{QR}=3ma$，隔离 QR，受力分析，$F_{PQ}=5ma$，联立解得 $k=\frac{3}{5}$，故 D 正确 C 错误。
}

\item
%% number: 10
%% paperName: 2017年海南理综第 10 题 5 分
%% typeId: 2
%% type: 多选题
%% chapter: 电磁感应
%% point: 线圈过有界磁场
%% method: 无
%% score: 5
%% degree: 450
%% duplicateId: 0
%% body:
如图，空间中存在一匀强磁场区域，磁场方向与竖直面（纸面）垂直，磁场的上、下边界（虚线）均为水平面；纸面内磁场上方有一个正方形导线框 abcd，其上、下两边均为磁场边界平行，边长小于磁场上、下边界的间距。若线框自由下落，从 ab 边进入磁场时开始，直至 ab 边到达磁场下边界为止，线框下落的速度大小可能\xzanswer{CD}
\onepicture[w=3.8cm]{figs/10.png}
\fourchoices[answer=CD]
{始终减小}
{始终不变}
{始终增加}
{先减小后增加}
%% answer: CD
%% memo:
\memoanswer{
若 ab 边进入磁场时所受安培力小于重力，则导线框继续加速运动，完全进入磁场后，做加速度为 $g$ 的加速运动，直到 ab 边到达磁场下边界，即线框下落的速度大小可能始终增加，故 C 正确；若 ab 边进入磁场时所受安培力大于重力，则导线框做减速运动，完全进入磁场后，做加速度为 $g$ 的加速运动，直到 ab 边到达磁场下边界，即线框下落的速度大小可能先减小后增加，故 D 正确。
}

\item
%% number: 11
%% paperName: 2017年海南理综第 11 题 6 分
%% typeId: 4
%% type: 实验题
%% chapter: 直线运动
%% point: 游标卡尺和螺旋测微仪
%% method: 无
%% score: 6
%% degree: 850
%% duplicateId: 0
%% body:
某同学用游标卡尺分别测量金属圆管的内、外壁直径，游标卡尺的示数分别如图（a）和图（b）所示。
\twopicture[align=center, wA=5.6cm, wB=5.6cm]
{figs/11a.png}
{figs/11b.png}

由图可读出，圆管内壁的直径为\tkanswer[1.6cm]{2.23}$\Ucm$，圆管外壁的直径为\tkanswer[1.6cm]{2.99}$\Ucm$；由此可计算出金属圆管横截面的面积。
%% answer:
\jdanswer{
$2.23\Ucm$，$2.99\Ucm$。
}
%% memo:
\memoanswer{
图 a 中游标卡尺的主尺读数为 $22\Umm$，游标尺上第 3 个刻度和主尺上某一刻度对齐，所以游标读数为 $3\times0.1\Umm=0.3\Umm$，所以最终读数为 $22\Umm+0.3\Umm=22.3\Umm=2.23\Ucm$。图 b 中游标卡尺的主尺读数为 $29\Umm$，游标尺上第 9 个刻度和主尺上某一刻度对齐，所以游标读数为 $9\times0.1\Umm=0.9\Umm$，所以最终读数为 $29\Umm+0.9\Umm=29.9\Umm=2.99\Ucm$。
}

\item
%% number: 12
%% paperName: 2017年海南理综第 12 题 12 分
%% typeId: 4
%% type: 实验题
%% chapter: 电路
%% point: 伏安法测电阻
%% method: 无
%% score: 12
%% degree: 550
%% duplicateId: 0
%% body:
某同学用伏安法测量待测电阻的阻值。现有器材为：待测电阻 $R$（阻值约为 $5\UO$），电源（电动势 $3\UV$），滑动变阻器（阻值范围 $0\sim10\UO$），电流表（量程 $0.6\UA$，$3\UA$），电压表（量程 $3\UV$，$15\UV$），开关，导线若干。实验要求在测量电路中将电流表外接，滑动变阻器起限流作用。回答下列问题：
\begin{enumerate}
	\item
	按照实验要求在图（a）中画出实物连线图。
	\onepicture[w=6.5cm, num=a]{figs/12a.png}
	\item
	若已按实验要求接线，闭合开关后移动滑动变阻器的滑片，电压表的示数始终约为 $3\UV$，电流表的示数始终接近 $0$。写出产生这种现象的一个原因：\tkanswer[3cm]{待测电阻 $R$ 断路}。
	\item
	在连线正确后，闭合开关。电压表和电流表的示数分别如图（b）和图（c）所示。由图可知，电压表读数为\tkanswer[1.4cm]{2.20}$\UV$，电流表读数为\tkanswer[1.4cm]{0.48}$\UA$。由此可得待测电阻的阻值为\tkanswer[1.8cm]{$4.58\UO$}（结果保留 3 位有效数字）。
	\twopicture[align=center, wA=5.6cm, wB=5.6cm, numA=b, numB=c]
	{figs/12b.png}
	{figs/12c.png}
\end{enumerate}
%% answer:
\jdanswer{
\begin{enumerate}
	\item
	如图所示。
	\item
	待测电阻 $R$ 断路。
	\item
	$2.20\UV$，$0.48\UA$，$4.58\UO$。
\end{enumerate}
}
%% memo:
\memoanswer{
（1）因为电源电动势为 $3\UV$，则电压表的量程选用 $3\UV$，根据欧姆定律知，电流的最大值大约 $0.6\UA$，则电流表量程选择 $0.6\UA$，根据实物图进行连线，实物连线图如答图所示。

\onepicture[w=6cm]{figs/12_an.png}

（2）闭合开关后移动滑动变阻器的滑片，电压表的示数始终约为 $3\UV$，电流表的示数始终接近 $0$，产生这种现象的原因是待测电阻 $R$ 断路，由于电压表内阻非常大，导致电流表电流接近 $0$，电压表电压测得是电源电压。

（3）由图可知，电压表的读数为 $2.20\UV$，电流表的读数为 $0.48\UA$，根据欧姆定律得，待测电阻 $R=\frac{U}{I}=\frac{2.20}{0.48}\approx4.58\UO$。
}

\item
%% number: 13
%% paperName: 2017年海南理综第 13 题 10 分
%% typeId: 6
%% type: 计算题
%% chapter: 电磁感应
%% point: 电磁感应受力分析
%% method: 无
%% score: 10
%% degree: 650
%% duplicateId: 0
%% body:
如图，两光滑平行金属导轨置于水平面（纸面）内，轨间距为 $l$，左端连有阻值为 $R$ 的电阻。一金属杆置于导轨上，金属杆右侧存在一磁感应强度大小为 $B$、方向竖直向下的匀强磁场区域。已知金属杆以速度 $v_{0}$ 向右进入磁场区域，做匀变速直线运动，到达磁场区域右边界（图中虚线位置）时速度恰好为零。金属杆与导轨始终保持垂直且接触良好。除左端所连电阻外，其他电阻忽略不计。求金属杆运动到磁场区域正中间时所受安培力的大小及此时电流的功率。
\onepicture[w=6cm, align=right]{figs/13.png}
%% answer:
\jdanswer{
$F=\frac{\sqrt{2}}{2R}B^{2}l^{2}v_{0}$，$P=\frac{1}{2R}B^{2}l^{2}v_{0}^{2}$。
}
%% memo:
\memoanswer{
设金属杆运动到磁场区域正中间时速度为 $v$，加速度大小为 $a$，位移为 $x$，由匀变速直线运动规律得 $v^{2}-v_{0}^{2}=-2ax$，又 $0-v_{0}^{2}=-2a(2x)$，联立解得 $v=\frac{\sqrt{2}}{2}v_{0}$，则此时产生的感应电动势为 $E=Blv=\frac{\sqrt{2}}{2}Blv_{0}$，金属杆中电流为 $I=\frac{E}{R}=\frac{\sqrt{2}}{2R}Blv_{0}$，金属杆所受安培力的大小 $F=BIl=\frac{\sqrt{2}}{2R}B^{2}l^{2}v_{0}$，此时电流的功率 $P=EI=\frac{1}{2R}B^{2}l^{2}v_{0}^{2}$。
}

\item
%% number: 14
%% paperName: 2017年海南理综第 14 题 16 分
%% typeId: 6
%% type: 计算题
%% chapter: 运动和力
%% point: 牛顿定律的应用
%% method: 无
%% score: 16
%% degree: 350
%% duplicateId: 0
%% body:
一轻弹簧的一端固定在倾角为 $\theta$ 的固定光滑斜面的底部，另一端和质量为 $m$ 的小物块 a 相连，如图所示。质量为 $\frac{3}{5}m$ 的小物块 b 紧靠 a 静止在斜面上，此时弹簧的压缩量为 $x_{0}$，从 $t=0$ 时开始，对 b 施加沿斜面向上的外力，使 b 始终做匀加速直线运动。经过一段时间后，物块 a、b 分离；再经过同样长的时间，b 距其出发点的距离恰好也为 $x_{0}$。弹簧的形变始终在弹性限度内，重力加速度大小为 $g$。求：
\begin{enumerate}
	\item
	弹簧的劲度系数；
	\item
	物块 b 加速度的大小；
	\item
	在物块 a、b 分离前，外力大小随时间变化的关系式。
\end{enumerate}
\onepicture[w=6.5cm, align=right]{figs/14.png}
%% answer:
\jdanswer{
\begin{enumerate}
	\item
	$\frac{8mg\sin\theta}{5x_{0}}$。
	\item
	$\frac{g\sin\theta}{5}$。
	\item
	$F=\frac{8}{25}mg\sin\theta+\frac{4mg^{2}\sin^{2}\theta}{25x_{0}}t^{2}$（$t<\sqrt{\frac{5x_{0}}{2g\sin\theta}}$）。
\end{enumerate}
}
%% memo:
\memoanswer{
（1）对整体分析，根据平衡条件可知，沿斜面方向上重力的分力与弹簧弹力平衡，则有 $kx_{0}=\left(m+\frac{3}{5}m\right)g\sin\theta$，解得 $k=\frac{8mg\sin\theta}{5x_{0}}$。

（2）由题意可知，b 经两段相等的时间位移为 $x_{0}$；由匀变速直线运动相邻相等时间内位移关系的规律可知 $\frac{x_{1}}{x_{0}}=\frac{1}{4}$，说明当形变量为 $x_{1}=x_{0}-\frac{x_{0}}{4}=\frac{3x_{0}}{4}$ 时二者分离；对 m 分析，因分离时 ab 间没有弹力，则根据牛顿第二定律可知 $kx_{1}-mg\sin\theta=ma$，联立解得 $a=\frac{g\sin\theta}{5}$。

（3）设时间为 $t$，则经时间 $t$ 时，ab 前进的位移 $x=\frac{1}{2}at^{2}=\frac{g\sin\theta t^{2}}{10}$，则形变量变为 $\Delta x=x_{0}-x$，对整体分析可知，由牛顿第二定律有 $F+k\Delta x-\left(m+\frac{3}{5}m\right)g\sin\theta=\left(m+\frac{3}{5}m\right)a$，解得 $F=\frac{8}{25}mg\sin\theta+\frac{4mg^{2}\sin^{2}\theta}{25x_{0}}t^{2}$。因分离时位移 $x=\frac{x_{0}}{4}$，由 $x=\frac{x_{0}}{4}=\frac{1}{2}at^{2}$ 解得 $t=\sqrt{\frac{5x_{0}}{2g\sin\theta}}$，故应保证 $t<\sqrt{\frac{5x_{0}}{2g\sin\theta}}$，$F$ 表达式才能成立。
}

\item
%% number: 15
%% paperName: 2017年海南理综第 15 题 8 分
%% typeId: 6
%% type: 计算题
%% chapter: 气体性质
%% point: 玻意耳定律
%% method: 无
%% score: 8
%% degree: 550
%% duplicateId: 0
%% body:
{}[选修 3-3]（12 分）
\begin{enumerate}
	\item
	（4 分）关于布朗运动，下列说法正确的是\tkanswer[1.6cm]{ABE}。（填入正确答案标号。选对 1 个得 2 分，选对 2 个得 3 分，选对 3 个得 4 分；有选错的得 0 分）
	\fivechoices[answer=ABE]
	{布朗运动是液体中悬浮微粒的无规则运动}
	{液体温度越高，液体中悬浮微粒的布朗运动越剧烈}
	{在液体中的悬浮颗粒只要大于某一尺寸，都会发生布朗运动}
	{液体中悬浮微粒的布朗运动使液体分子永不停息地做无规则运动}
	{液体中悬浮微粒的布朗运动是液体分子对它的撞击作用不平衡所引起的}
	\item
	（8 分）一粗细均匀的 U 形管 ABCD 的 A 端封闭，D 端与大气相通。用水银将一定质量的理想气体封闭在 U 形管的 AB 一侧，并将两端向下竖直放置，如图所示。此时 AB 侧的气体柱长度 $l_{1}=25\Ucm$。管中 AB、CD 两侧的水银面高度差 $h_{1}=5\Ucm$。现将 U 形管缓慢旋转 $180^{\circ}$，使 A、D 两端在上，在转动过程中没有水银漏出。已知大气压强 $p_{0}=76\UcmHg$。求旋转后，AB、CD 两侧的水银面高度差。
	\onepicture[w=3cm, align=right]{figs/15.png}
\end{enumerate}
%% answer:
\jdanswer{
\begin{enumerate}
	\item
	ABE。
	\item
	$1\Ucm$。
\end{enumerate}
}
%% memo:
\memoanswer{
（1）布朗运动是液体中悬浮微粒的无规则运动，液体温度越高，液体中悬浮微粒的布朗运动越剧烈，故 AB 正确；在液体中的悬浮颗粒只要小于某一尺寸，都会发生布朗运动，液体中悬浮微粒的布朗运动是液体分子永不停息地做无规则运动的反映，液体中悬浮微粒的布朗运动是液体分子对它的撞击作用不平衡所引起的，故 E 正确 CD 错误。

（2）设玻璃管横截面积为 $S$，初状态气体压强为 $p_{1}=p_{0}+h_{1}=76\UcmHg+5\UcmHg=81\UcmHg$①，气柱体积为 $V_{1}=l_{1}S=25S$②。

旋转 $180^{\circ}$ 后，设 AB 管中气柱长度增加 $\Delta h$，假设 AB 管中液面高于 CD 管中液面，即 $\Delta h<\frac{h_{1}}{2}=2.5\Ucm$。由几何知识可知两侧液面高度差 $h_{2}=h_{1}-2\Delta h=5-2\Delta h$③，气体压强为 $p_{2}=p_{0}-h_{2}$④，气柱体积为 $V_{2}=(l_{1}+\Delta h)S=(25+\Delta h)S$⑤。

温度一定，由玻意耳定律得 $p_{1}V_{1}=p_{2}V_{2}$⑥，联立①～⑥式得 $(p_{0}+h_{1})l_{1}S=(p_{0}-h_{1}+2\Delta h)(l_{1}+\Delta h)S$⑦。

将题中所给数据代入⑦式，解方程可得 $\Delta h=2\Ucm$ 或 $\Delta h=-62.5\Ucm$（不合题意舍去），则此时两管高度差 $h_{2}=h_{1}-2\Delta h=1\Ucm$，且 AB 管中液面高于 CD 管。
}

\item
%% number: 16
%% paperName: 2017年海南理综第 16 题 8 分
%% typeId: 6
%% type: 计算题
%% chapter: 机械波
%% point: 波的图象
%% method: 无
%% score: 8
%% degree: 450
%% duplicateId: 0
%% body:
{}[选修 3-4]（12 分）
\begin{enumerate}
	\item
	（4 分）如图，空气中有两块材质不同、上下表面平行的透明玻璃板平行放置；一细光束从空气中以某一角度 $\theta$（$0<\theta<90^{\circ}$）入射到第一块玻璃板的上表面。下列说法正确的是\tkanswer[1.6cm]{ACD}。（填入正确答案标号。选对 1 个得 2 分，选对 2 个得 3 分，选对 3 个得 4 分；有选错的得 0 分）
	\onepicture[w=5.5cm]{figs/16a.png}
	\fivechoices[answer=ACD]
	{在第一块玻璃板下表面一定有出射光}
	{在第二块玻璃板下表面一定没有出射光}
	{第二块玻璃板下表面的出射光方向一定与入射光方向平行}
	{第二块玻璃板下表面的出射光一定在入射光延长线的左侧}
	{第一块玻璃板下表面的出射光线一定在入射光延长线的右侧}
	\item
	（8 分）从两个波源发出的两列振幅相同、频率均为 $5\UHz$ 的简谐横波，分别沿 $x$ 轴正、负方向传播，在某一时刻到达 A、B 点，如图中实线、虚线所示。两列波的波速均为 $10\Ums$。求：
	\begin{enumerate}[label=(\roman*)]
		\item
		质点 P、O 开始振动的时刻之差；
		\item
		再经过半个周期后，两列波在 $x=1\Um$ 和 $x=5\Um$ 之间引起的合振动振幅极大和极小的质点的 $x$ 坐标。
	\end{enumerate}
	\onepicture[w=6.5cm, align=right]{figs/16b.png}
\end{enumerate}
%% answer:
\jdanswer{
\begin{enumerate}
	\item
	ACD。
	\item
	（i）$0.05\Us$；（ii）合振动振幅极大的质点的 $x$ 坐标为 $2\Um$、$3\Um$、$4\Um$，合振动振幅极小的质点的 $x$ 坐标为 $1.5\Um$、$2.5\Um$、$3.5\Um$、$4.5\Um$。
\end{enumerate}
}
%% memo:
\memoanswer{
（1）由折射定律知，光由空气斜射入玻璃，在两块玻璃板下表面一定有出射光，故 A 正确 B 错误；第二块玻璃板下表面的出射光方向一定与入射光方向平行，且出射光一定在入射光延长线的左侧，故 CD 正确 E 错误。

（2）（i）由题意知，两波周期均为 $T=\frac{1}{f}=0.2\Us$，由波形图可知，质点 P 在该时刻前 $0.35\Us$ 开始振动，质点 O 在该时刻前 $0.4\Us$ 开始振动，故 P、O 开始振动的时刻之差为 $0.05\Us$。

（ii）波长 $\lambda=\frac{v}{f}=2\Um$，画出再经过半个周期后，两列波的波形图，如图所示。

\onepicture[w=6cm]{figs/16_an.png}

由图可知，两列波在 $x=1\Um$ 和 $x=5\Um$ 之间引起的合振动振幅极大的质点的 $x$ 坐标为 $x=2\Um$、$x=3\Um$、$x=4\Um$，振幅极小的质点的 $x$ 坐标为 $x=1.5\Um$、$x=2.5\Um$、$x=3.5\Um$、$x=4.5\Um$。
}

\end{enumerate}

\end{document}
